\section{Discussion}
\label{sec:limits}
R20 produced six complete task-to-approval-to-receipt-to-sample traces across
two simulator scenes; none of its single post-readback samples reached target.
R21 passed 31 assertions on the same configured route, and a 20-case mock
suite refused twelve stale approval-context mutations across three recipe
actuation routes before actuation. A test-only ablation that removed the
post-approval tuple comparison allowed the product and active-run mutations
to proceed on all three routes, while independent checks still blocked line
and recipe changes. This isolates the added coverage of that comparison for
those two context fields; it is not a baseline comparison of AgentWorkShop
against another integration framework.

Each scene used one simulator instance and a single post-readback probe per
task; the results therefore do not establish independent process replication,
convergence, or quality improvement. Configuration reuse, alternate harnesses
or direct REST entry, online DAQ, physical transfer, and comparative
architecture value remain untested. Synthetic approval does not measure
operator effectiveness
\cite{parasuramanmanzey2010complacency,parasuraman2000model,bainbridge1983ironies};
tool restrictions do not establish deployment security
\cite{greshake2023injection,liu2024pinjection,iec62443-3-3} or replace
independently engineered PLC/SIS protection \cite{iec61511}. Future work
should compare a defined platform-level baseline and include independent
simulator and physical-plant trials.

\section{Conclusion}
\label{sec:conclusion}
AgentWorkShop connects collaborative tasks to production context, checked
operations, semantic I/O, and service-owned records. R20 and R21 verify
selected paths across two simulator mappings; a mock-state suite confirms
refusal when an approved line--product--recipe--run tuple changes on three
actuation routes. A targeted ablation shows that the tuple comparison adds
mock-state refusal coverage for product and active-run changes beyond the
remaining route checks. These results validate selected software behavior and
leave platform-level comparative integration value and process performance
for future evaluation.
